% =========================================================
%  sesi/04/sesi.tex
%  Sesi 4 — Gravitasi, lompatan, dan pantulan
% =========================================================

\sesi{04}{Gravitasi, Lompatan, dan Pantulan}{1.3}{130 menit}

Sesi ini adalah pertemuan minggu ke-4. Sejak sesi sebelumnya pemain sudah
dapat berjalan ke kiri dan ke kanan dengan kecepatan 260 piksel per detik,
tetapi geraknya hanya mendatar: tidak ada yang menariknya ke bawah, dan tidak
ada yang menghalanginya selain tepi jendela. Gudang Artefak masih terasa
seperti lorong, bukan gudang.

Pada sesi ini Anda menambahkan arah tegak. Pemain jatuh karena
\istilah{gravitasi}, melompat dengan \istilah{jump impulse} ketika kakinya
menyentuh lantai atau rak, dan memantul ketika membentur dinding. Lantai serta
rak berubah dari gambar menjadi batas yang benar-benar dihitung, sehingga
pemain dapat berdiri di atasnya.

Seluruh perilaku itu dihasilkan oleh tiga baris perhitungan yang sama untuk
setiap benda. Angka yang dipakai tetap: gravitasi 1900 piksel per detik
kuadrat, kecepatan lompatan 720 piksel per detik, \istilah{koefisien
restitusi} dinding 0,8, dan ambang pantulan 60 piksel per detik. Sesudah sesi
ini, kerangka gerak pemain sudah lengkap untuk sebuah permainan platform
sederhana.

% ---------------------------------------------------------
\bagian{Tujuan Sesi}

\begin{tujuanbox}
Setelah menuntaskan sesi ini, Anda dapat menunjukkan:
\begin{enumerate}[leftmargin=1.5em, itemsep=0.2em]
  \item pemain yang jatuh dengan percepatan 1900 piksel per detik kuadrat dan
        berhenti tepat di atas lantai atau rak;
  \item lompatan setinggi sekitar 136 piksel yang hanya terjadi ketika pemain
        menyentuh permukaan, dan tidak berulang walaupun tombol ditahan;
  \item pantulan dinding dengan koefisien restitusi 0,8 yang berhenti sendiri
        setelah kecepatannya di bawah 60 piksel per detik;
  \item penjelasan urutan pembaruan \istilah{integrasi Euler}: kecepatan
        diperbarui lebih dahulu, posisi menyusul;
  \item pembacaan nilai \code{GRAVITY}, \code{JUMP_SPEED},
        \code{WALL_RESTITUTION}, dan \code{MIN_BOUNCE_SPEED} pada
        \file{settings.py}, beserta akibat yang muncul ketika salah satunya
        diubah.
\end{enumerate}
\end{tujuanbox}

% ---------------------------------------------------------
\bagian{Prasyarat}

\begin{prasyaratbox}
\begin{itemize}[leftmargin=1.4em, itemsep=0.25em]
  \item Anda sudah menjalankan checkpoint sesi 3: pemain bergerak mendatar
        mengikuti tombol, dan posisinya disimpan sebagai \code{Vector2}
        pecahan.
  \item Anda memahami pembagian tugas antara \code{position} yang bertipe
        pecahan dan \code{rect} yang bertipe bilangan bulat.
  \item Anda dapat menjelaskan mengapa kecepatan dikalikan
        \code{delta_seconds} sebelum ditambahkan ke posisi.
  \item Anda dapat menjalankan program dengan pilihan \code{--frames} untuk
        menghentikannya secara otomatis.
\end{itemize}
\end{prasyaratbox}

% ---------------------------------------------------------
\bagian{Peta Kemampuan}

\petasesi
  {Pemain bergerak mendatar dengan kecepatan tetap; posisi pecahan
   diselaraskan dengan \code{rect}.}
  {Gravitasi 1900, lompatan 720, pantulan dinding 0,8, ambang 60, serta lantai
   dan rak sebagai batas gerak.}
  {Kelas \code{Sprite} (\istilah{sprite}), wadah \code{Group}
   (\istilah{group}), dan \istilah{collision} dengan \istilah{AABB} pada
   sesi 5.}

% ---------------------------------------------------------
\bagian{Konsep yang Dipakai}

\subsection*{Percepatan mengubah kecepatan, kecepatan mengubah posisi}

Sampai sesi 3, kecepatan pemain hanya berubah ketika tombol ditekan atau
dilepas. Pada sesi ini kecepatan tegak berubah dengan sendirinya setiap
bingkai. Perubahan itu dihitung dengan dua langkah yang berurutan:

\[
v \leftarrow v + a\,\Delta t
\qquad\text{lalu}\qquad
p \leftarrow p + v\,\Delta t .
\]

Di sini $v$ adalah kecepatan tegak dalam piksel per detik, $a$ adalah
percepatan gravitasi dalam piksel per detik kuadrat, $\Delta t$ adalah selang
waktu satu bingkai dalam detik, dan $p$ adalah posisi tegak dalam piksel.
Karena sumbu tegak layar mengarah ke bawah, gravitasi yang menarik benda ke
bawah diberi tanda positif, sedangkan lompatan yang mengarah ke atas diberi
tanda negatif.

Urutan kedua baris itu penting dan disebut \istilah{integrasi Euler}
semi-implisit. Posisi memakai kecepatan yang \emph{baru saja} diperbarui,
bukan kecepatan bingkai sebelumnya. Apabila urutannya dibalik, perpindahan
selalu tertinggal satu bingkai; pada satu bingkai selisihnya kecil, tetapi
selama ratusan bingkai ia menumpuk dan membuat tinggi lompatan serta titik
mendarat meleset dari perhitungan.

\subsection*{Lompatan adalah impuls, bukan gaya}

\istilah{Jump impulse} adalah perubahan kecepatan yang terjadi seketika.
Karena itu nilainya \emph{diganti}, bukan ditambahkan setiap bingkai:
\code{velocity.y = -JUMP_SPEED}. Tanda negatif membuat pemain bergerak ke
atas. Sesudah itu gravitasi bekerja seperti biasa, memperlambat pemain sampai
berhenti di puncak, lalu menariknya kembali ke bawah.

Lompatan hanya boleh terjadi ketika \code{on_ground} bernilai benar. Tanpa
syarat itu pemain dapat melompat di udara berkali-kali. Karena lompatan adalah
peristiwa sekali, tombolnya dibaca dari \code{KEYDOWN}, bukan dari keadaan
tombol yang terus benar selama tombol ditahan.

\subsection*{Mendarat, memantul, dan ambang}

Ketika badan pemain menyentuh permukaan, posisinya dikembalikan ke tepi
permukaan itu lebih dahulu, baru kecepatannya diubah. Untuk dinding, arah
gerak dibalik dan besarnya dikurangi dengan \istilah{koefisien restitusi} $e$:

\[
v' = -e\,v .
\]

Dengan $e = 0{,}8$, kecepatan sesudah pantulan tinggal 80 persen kecepatan
datangnya. Pantulan dihentikan ketika kecepatannya sudah di bawah 60 piksel
per detik, sebab di bawah nilai itu pantulan hanya mengangkat pemain kurang
dari satu piksel dan tidak terlihat, tetapi tetap dihitung setiap bingkai.
Pada sumbu tegak, pemain tidak memantul: kecepatannya langsung dinolkan agar
mendarat dengan mantap.

\subsection*{Satu sumbu pada satu waktu}

Gerak mendatar dan gerak tegak diperiksa terpisah, mendatar lebih dahulu.
Sesudah sumbu mendatar selesai dan tidak ada lagi tumpang tindih, barulah
sumbu tegak diperiksa. Cara ini membuat arah datang dapat disimpulkan dari
tanda kecepatan: bergerak ke kanan berarti menabrak sisi kiri benda, dan
sebaliknya. Bila kedua sumbu diperiksa sekaligus, program hanya tahu bahwa
kedua persegi beririsan, lalu sering memilih penyelesaian yang salah.

% ---------------------------------------------------------
\bagian{Tutorial}

Kerjakan langkah berikut berurutan. Setiap langkah diakhiri oleh perubahan
yang dapat Anda jalankan dan lihat akibatnya.

\begin{langkah}
  \item \textbf{Tambahkan konstanta fisika.} Buka
        \file{gudang_artefak/settings.py}, lalu tambahkan \code{GRAVITY =
        1900.0}, \code{JUMP_SPEED = 720.0}, \code{WALL_RESTITUTION = 0.8},
        dan \code{MIN_BOUNCE_SPEED = 60.0} pada bagian pemain dan fisika.
        Tulis satuannya sebagai komentar agar angka itu tidak menjadi teka-teki
        di kemudian hari.

  \item \textbf{Lengkapi keadaan pemain.} Di dalam \code{Player.__init__},
        tambahkan \code{self.velocity = pygame.Vector2(0, 0)} dan
        \code{self.on_ground = False}. Ukuran pemain tetap 36 piksel seperti
        pada sesi 3, sehingga \code{rect} dan \code{position} tetap mengacu
        pada badan yang sama. Posisi mendatar yang sudah ada tetap
        dipertahankan; keduanya kini menjadi sumber kebenaran gerak.

  \item \textbf{Tulis pembaruan gerak.} Ganti metode \code{update} menjadi
        \code{update_side} yang menerima \code{delta_seconds}, arah tombol, dan
        daftar permukaan padat. Isinya: tetapkan kecepatan mendatar, tambahkan
        gravitasi ke kecepatan tegak, lalu kosongkan \code{on_ground} sebelum
        pemeriksaan ulang.

  \item \textbf{Siapkan lantai dan rak.} Pada \file{gudang_artefak/levels.py},
        tulis \code{PLATFORM_LAYOUT} berisi enam persegi panjang: lantai
        setinggi 48 piksel di tepi bawah, dan lima rak tipis. Fungsi
        \code{build_rak()} mengembalikan daftarnya sebagai \code{pygame.Rect}.
        Simpan hasilnya di \code{Game} sebagai \code{self.solids}.

  \item \textbf{Selesaikan sumbu mendatar.} Pindahkan posisi pada sumbu $x$,
        selaraskan \code{rect.x} dengan \code{round()}, lalu telusuri
        \code{solids}. Bila \code{colliderect} benar, tempelkan sisi pemain ke
        sisi benda, lalu balik kecepatan mendatar dengan koefisien 0,8 atau
        nolkan bila besarnya di bawah 60.

  \item \textbf{Selesaikan sumbu tegak.} Pindahkan posisi pada sumbu $y$,
        selaraskan \code{rect.y}, lalu periksa kembali seluruh \code{solids}.
        Bila pemain sedang bergerak ke bawah, letakkan \code{rect.bottom} di
        \code{solid.top}, nolkan kecepatan tegak, dan tandai
        \code{on_ground = True}. Bila sedang bergerak ke atas, letakkan
        \code{rect.top} di \code{solid.bottom} dan nolkan kecepatannya.

  \item \textbf{Tambahkan lompatan.} Tulis metode \code{jump()} yang hanya
        bekerja bila \code{on_ground} benar, lalu panggil dari peristiwa
        \code{KEYDOWN} untuk tombol spasi. Jangan memakai
        \code{key.get_pressed()} untuk lompatan.

  \item \textbf{Jalankan dan uji.} Gambar seluruh \code{solids} dengan
        \code{pygame.draw.rect}, lalu jalankan program. Lompat dari lantai ke
        rak pertama, dan coba tahan tombol spasi untuk memastikan pemain tidak
        terbang.
\end{langkah}

\begin{ujibox}
Setelah Langkah 8, periksa lima hal berikut.
\begin{itemize}[leftmargin=1.4em, itemsep=0.2em]
  \item Pemain jatuh dari titik awal dan berhenti dengan kakinya tepat di
        permukaan lantai, tidak terbenam dan tidak melayang.
  \item Lompatan dari lantai mencapai tinggi sekitar 136 piksel, cukup untuk
        rak pertama pada ketinggian 114 piksel di atas lantai.
  \item Menahan spasi hanya menghasilkan satu lompatan. Bila pemain terus
        melompat, periksa sumber peristiwanya.
  \item Berjalan ke dinding membuat pemain memantul sekali atau dua kali,
        lalu berhenti; pemain tidak bergetar di dinding.
  \item Program berhenti rapi dengan \code{--frames 120} pada laju 30 maupun
        60 bingkai per detik, dan tinggi lompatannya tidak berbeda jauh.
\end{itemize}
\end{ujibox}

% ---------------------------------------------------------
\bagian{Kode Sesi Ini}

Potongan berikut adalah bagian yang berubah pada checkpoint sesi ini. Berkas
\file{levels.py} memuat susunan lantai dan rak, sedangkan \file{entities.py}
memuat seluruh perhitungan gerak. Bacalah dengan urutan dari atas ke bawah;
setiap potongan bergantung pada potongan sebelumnya.

\tampilkodebaris[Python]{code/sesi-04/gudang_artefak/settings.py}{22}{28}
  {settings.py: konstanta pemain dan fisika}

\tampilkodebaris[Python]{code/sesi-04/gudang_artefak/entities.py}{27}{47}
  {entities.py: pembaruan gerak, urutan Euler, dan dua sumbu}

\tampilkodebaris[Python]{code/sesi-04/gudang_artefak/entities.py}{49}{63}
  {entities.py: penyelesaian sumbu mendatar dan pantulan dinding}

\tampilkodebaris[Python]{code/sesi-04/gudang_artefak/entities.py}{65}{76}
  {entities.py: penyelesaian sumbu tegak dan penanda menyentuh permukaan}

\tampilkodebaris[Python]{code/sesi-04/gudang_artefak/entities.py}{78}{83}
  {entities.py: lompatan yang hanya bekerja di atas permukaan}

\tampilkodebaris[Python]{code/sesi-04/gudang_artefak/levels.py}{12}{29}
  {levels.py: susunan lantai dan rak sebagai Rect}

\tampilkodebaris[Python]{code/sesi-04/gudang_artefak/game.py}{56}{72}
  {game.py: peristiwa lompatan dan pembaruan fisika}

\checkpoint{04}{Pemain dapat berjalan, jatuh, mendarat di lantai atau rak,
melompat sekali setiap penekanan spasi, dan memantul di dinding dengan
kecepatan yang makin mengecil.}

% ---------------------------------------------------------
\bagian{Kesalahan Umum}

\begin{galatbox}
\begin{itemize}[leftmargin=1.4em, itemsep=0.35em]
  \item \textbf{Pemain melayang ke atas, bukan jatuh.}
        Gejala: pemain perlahan naik dan hilang di tepi atas jendela. Periksa
        tanda \code{GRAVITY}; pada koordinat layar nilai positif menarik ke
        bawah. Bila Anda menulisnya negatif, gravitasi bekerja seperti
        dorongan ke atas.
  \item \textbf{Pemain melompat terus selama spasi ditahan.}
        Gejala: pemain terbang seperti roket. Periksa bahwa lompatan dibaca
        dari \code{KEYDOWN}, bukan dari \code{pygame.key.get_pressed()}.
        Peristiwa hanya muncul sekali pada setiap penekanan.
  \item \textbf{Pemain bergetar di permukaan lantai.}
        Gejala: pemain tampak bergetar halus dan \code{on_ground} berganti
        nilai setiap bingkai. Periksa apakah kecepatan tegak dinolkan sesudah
        posisi diperbaiki. Tanpa penolan itu, gravitasi menekan pemain kembali
        ke dalam lantai pada bingkai berikutnya.
  \item \textbf{Pemain menembus rak yang tipis.}
        Gejala: pemain melewati rak tanpa pernah mendarat. Periksa pemeriksaan
        tabrakan; jangan memakai tanda kesamaan (\code{==}) pada posisi.
        Gunakan \code{colliderect} dan lakukan pemeriksaan sesudah posisi
        diperbarui. Bila kecepatannya terlalu besar, bagi perpindahan atau
        batasi kecepatan jatuh.
  \item \textbf{Posisi dan \code{rect} tidak lagi sama.}
        Gejala: pemain mendarat di satu tempat, lalu melompat dari tempat lain
        yang beberapa piksel berbeda. Periksa bahwa sesudah setiap perbaikan
        \code{rect}, nilai \code{self.position} disalin kembali dari
        \code{rect}.
\end{itemize}
\end{galatbox}

% ---------------------------------------------------------
\bagian{Lembar Kerja}

Isi tabel berikut pada catatan praktikum Anda. Dua baris pertama meminta
perbandingan antara hasil pengukuran dan perhitungan, bukan sekadar angka
pengamatan.

{\small
\renewcommand{\arraystretch}{1.4}
\begin{tabularx}{\linewidth}{@{}C{0.8cm} L{5.4cm} Y@{}}
\toprule
\textbf{No} & \textbf{Hasil yang dicatat} & \textbf{Keterangan Anda} \\
\midrule
1 & Tinggi lompatan terukur dan hasil $v_0^2/(2g)$ & \\[0.55cm]
2 & Lama pemain di udara dan hasil $2v_0/g$ & \\[0.55cm]
3 & Kecepatan pantulan dinding pertama dan kedua & \\[0.55cm]
4 & Nilai \code{MIN_BOUNCE_SPEED} sebelum dan sesudah diubah & \\[0.55cm]
5 & Satu kendala yang Anda hadapi dan cara menyelesaikannya & \\[0.9cm]
\bottomrule
\end{tabularx}
}

Dengan $v_0 = 720$ piksel per detik dan $g = 1900$ piksel per detik kuadrat,
hitungan memberi tinggi sekitar 136 piksel dan lama di udara sekitar 0,76
detik. Selisih beberapa piksel dari pengukuran adalah galat integrasi Euler
yang wajar, bukan kesalahan program.

% ---------------------------------------------------------
\bagian{Checklist Sesi}

\begin{checklistbox}
\begin{checklist}
  \item Saya dapat menunjukkan pemain yang jatuh dan berhenti di atas lantai
        serta rak.
  \item Saya dapat menjelaskan arti setiap lambang pada kedua rumus pembaruan
        gerak.
  \item Saya dapat menjelaskan mengapa kecepatan diperbarui sebelum posisi.
  \item Saya dapat menunjukkan lompatan yang hanya terjadi ketika pemain
        menyentuh permukaan.
  \item Saya dapat menunjukkan pantulan dinding 0,8 yang berhenti sendiri di
        bawah ambang 60 piksel per detik.
  \item Saya dapat mengubah satu konstanta fisika dan menerangkan akibat
        perubahan itu pada gerak pemain.
\end{checklist}
\end{checklistbox}

% ---------------------------------------------------------
\bagian{Pengait ke Sesi Berikutnya}

Fisika pada sesi ini sudah benar, tetapi seluruh isi gudang masih berupa
daftar persegi panjang dan pemain masih satu-satunya benda yang hidup. Pada
sesi 5 Anda mengubah lantai, rak, artefak, peti, dan pintu menjadi
\istilah{sprite} yang dikelola \istilah{group}, lalu memakai \istilah{AABB}
untuk memeriksa persentuhannya. Di sanalah permainan mulai memiliki aturan:
artefak menambah skor, peti mengurangi nyawa, dan pintu terbuka setelah
seluruh artefak terkumpul.
